% TOP_PROBLEM: {"id": 30006744, "title": "NP-Completeness of the Minimum Circuit Size Problem (MCSP)", "category_id": 15, "category": {"id": 15, "name": "computer_science", "display_name": "Computer Science", "description": "Computational complexity, algorithms, and theoretical CS.", "slug": "computer-science", "order_index": 15, "created_at": "2026-07-31T15:26:25.671Z"}, "status": "open"}
% FIRST_POSTED: null
% !TeX program = lualatex
% Compile twice: lualatex 30006744.tex
% All references and prose are embedded; no BibTeX database or external inputs.
\documentclass[11pt,a4paper]{article}
\usepackage[margin=24mm,headheight=14pt]{geometry}
\usepackage{fontspec}
\setmainfont{Latin Modern Roman}
\setsansfont{Latin Modern Sans}
\setmonofont{Latin Modern Mono}
\usepackage{amsmath,amssymb,mathtools}
\usepackage{unicode-math}
\setmathfont{Latin Modern Math}
\usepackage{microtype}
\usepackage{enumitem,xurl,needspace}
\usepackage[dvipsnames]{xcolor}
\usepackage{hyperref}
\hypersetup{unicode=true,colorlinks=true,linkcolor=MidnightBlue,urlcolor=MidnightBlue,
 pdftitle={Research attempt: Problem 30006744},pdfauthor={Research notebook},bookmarksnumbered=true}
\usepackage{bookmark}
\usepackage{tocloft}
\setlength{\cftsecnumwidth}{3em}
\cftsetpnumwidth{2.5em}
\cftsetrmarg{4em}
\usepackage{titlesec}
\titleformat{\section}{\Large\bfseries\raggedright}{\thesection}{0.7em}{}
\usepackage{fancyhdr}
\pagestyle{fancy}\fancyhf{}
\fancyhead[L]{\small\sffamily Open Problems Math | Research notebook}
\fancyhead[R]{\small Problem 30006744 | 27 September 2026}
\fancyfoot[C]{\thepage}
\setlength{\parindent}{0pt}
\setlength{\parskip}{3pt plus 1pt}
\setlength{\emergencystretch}{3em}
\setcounter{tocdepth}{1}
\setcounter{secnumdepth}{1}
\setlist[itemize]{leftmargin=3em,itemsep=5pt,topsep=4pt,parsep=0pt}
\allowdisplaybreaks
\newcommand{\N}{\mathbb{N}}
\newcommand{\Z}{\mathbb{Z}}
\newcommand{\Q}{\mathbb{Q}}
\newcommand{\R}{\mathbb{R}}
\newcommand{\C}{\mathbb{C}}
\newcommand{\F}{\mathbb{F}}
\newcommand{\A}{\mathbb{A}}
\newcommand{\PP}{\mathbb P}
\newcommand{\E}{\mathbb{E}}
\newcommand{\eps}{\varepsilon}
\newcommand{\dd}{\,\mathrm d}
\newcommand{\sref}[2]{\hyperlink{source:#1:#2}{[S#2]}}
\newcommand{\eref}[2]{\hyperlink{extra:#1:#2}{[E#2]}}
% Turn the next switch off for a shorter reading copy. The quotations preserve
% every exactTarget field, including qualifications omitted from a short summary.
\newif\ifshowcatalogue
\showcataloguetrue
\newcommand{\cataloguescope}[1]{\ifshowcatalogue
 \paragraph{Exact scope in the supplied catalogue.}
 \begin{quote}\small #1\end{quote}\fi}
\usepackage{etoolbox}
\AtBeginEnvironment{itemize}{\small}
\numberwithin{equation}{section}
% Standalone export. Original research content preserved.
\begin{document}
\setcounter{section}{0}
% ================================================================
% problemId: 30006744
% Canonical problem ID: 30006744
% exactTarget SHA-256: 510c09a28f82a4eb055322b0840aadddb6d05e6783f78c1d9a57b6e8073ad94e
\clearpage
\section*{\texorpdfstring{NP-Completeness of the Minimum Circuit Size Problem (MCSP)}{NP-Completeness of the Minimum Circuit Size Problem (MCSP)}}\label{30006744}
\subsection{Definitions and mathematical statement}
The input to MCSP is the full $2^n$-bit truth table of $f:\{0,1\}^n\to\{0,1\}$ and an integer size bound $s$. Using a fixed standard complete Boolean basis, define
\[
 \mathsf{MCSP}=\{(\operatorname{tt}(f),s):\exists C,
                          |C|\le s,\ C(x)=f(x)\ \forall x\in\{0,1\}^n\}.
\]
The question is NP-hardness under deterministic polynomial-time many-one reductions, with running time measured in truth-table input length. A succinct or implicitly specified truth table defines a different problem. Randomized or non-many-one hardness reductions do not establish the selected reduction type.
\par\smallskip\textit{Formulation and concept references:} \sref{30006744}{1}.
\cataloguescope{Prove or disprove that MCSP is NP-hard under deterministic polynomial-time many-one reductions, where an input is the truth table of a Boolean function and a size bound and asks whether a circuit of at most that size computes the function.}
\subsection{Short English statement}
Is finding whether a complete Boolean truth table has a small circuit NP-hard under ordinary deterministic many-one reductions?
% BEGIN RESEARCH ADDITION 134
\subsection{Research attempt}

\paragraph{Current frontier.}
The June 2026 formulation source explicitly leaves open deterministic,
polynomial-time, many-one hardness for the \emph{explicit} truth-table
problem. Its new hardness result concerns an implicitly sampled function,
uses randomness, and assumes subexponential cryptography and proof-complexity
hypotheses. None of those qualifications can be discarded.
\sref{30006744}{1} Hirahara--Ilango obtain conditional hardness with
quasipolynomial-time nonadaptive reductions, not the reduction requested
here. \eref{30006744}{1}

\paragraph{Attack route.}
Try to transfer hardness from minimization of partial truth tables to total
ones. A partial table $p:\{0,1\}^n\to\{0,1,*\}$ specifies only some values;
write $C^*(p)=\min_{f\supset p}C(f)$. Three possible routes are a
canonical completion, a selector encoding all completions, and a robust gap
gadget. The first two can be tested without assuming a conjecture. The
third needs a totalization theorem with a controlled gap for unrestricted
circuits, not just a gadget working on the intended representations.
The following elementary deductions are developed for this notebook;
no historical novelty is asserted.

\paragraph{Attempt: canonical completion loses the existential witness.}
Use first the basis of all binary gates, free input wires and free constants;
a fixed complete basis changes the upper bounds below by constant factors.
For any total $h$, put
\[
 p_h(x)=\begin{cases}1&h(x)=1,\\ *&h(x)=0.\end{cases}
\]
Then $C^*(p_h)=0$, using constant one, whereas zero-filling produces exactly
$h$. Thus zero-filling can turn a trivial partial instance into an arbitrarily
hard total function. One-filling has the symmetric defect using specified
zeros. Charging one gate for a constant only changes a fixed additive cost.

For completeness, circuits with $s$ gates on $n$ inputs have at most
\[
 (n+s+2)\bigl(16(n+s+2)^2\bigr)^s
\]
ordered descriptions: choose predecessors, a gate truth table, and an output
wire. Circuits with fewer gates can be padded by projection gates. For fixed
$s$ this is less than $2^{2^n}$ for sufficiently large $n$, so some $h$ has
$C(h)>s$. This counting argument does not efficiently identify a hard table
at the scale a reduction would need.

Could an address register retain the choice? If $p$ has $u$ stars, a table
$F(z,x)$ displaying the completion selected by $z\in\{0,1\}^u$ has
$2^{n+u}$ entries. Relative to input length $N=2^n$, polynomial output length
requires $u=O(\log N)$, not $u=\Theta(N)$. Moreover restriction proves
\[
 C(F)\ge\max_z C(F|_z),
\]
not the minimum over $z$. Computing every completion reverses the
existential quantifier in $C^*(p)$.

\paragraph{An edit-distance estimate for future gap gadgets.}
Suppose $f,g$ differ on exactly $k\ge1$ rows $E$. For each $a\in E$ compute
\[
 e_a(x)=\bigwedge_{i=1}^n
 \begin{cases}x_i&a_i=1,\\ \neg x_i&a_i=0.
 \end{cases}
\]
Share the $n$ negated inputs; each conjunction uses $n-1$ gates, the disjoint
indicators can be joined by $k-1$ OR gates, and
$g=f\mathbin{\oplus}\bigvee_{a\in E}e_a$ uses one final binary gate.
Reversing the roles of $f,g$ proves
\begin{equation}\label{30006744:edit}
 |C(f)-C(g)|\le n+k(n-1)+(k-1)+1=n(k+1).
\end{equation}
For $k=0$ the difference is zero. A fixed complete basis gives the bound
$c_{\mathcal B}n(k+1)$ by a fixed gate simulation.

Consequently, if a proposed reduction has a proved yes-bound $a(m)$ and
no-bound $b(m)$ before a correction of at most $k(m)$ table entries, its gap
survives whenever
\[
 b(m)-a(m)>2c_{\mathcal B}n(m)(k(m)+1).
\]
The corrected bounds are respectively $a+c_{\mathcal B}n(k+1)$ and
$b-c_{\mathcal B}n(k+1)$; an integer threshold strictly below the latter
and at least the former works. This gives a concrete error budget, but
neither the initial gap nor a suitable correction algorithm. On $p_h$, the
complexity of the unspecified mask is exactly the uncontrolled quantity.

Repetition supplies no automatic amplification. If $F(z,x)=f(x)$, ignoring
$z$ and restricting it give $C(F)=C(f)$ in the auxiliary convention.
Repeating the same table expands output length but does not multiply
complexity. XOR of two identical slices is constant.

\paragraph{Stress test and dependency check.}
These are failures of specified transformations, not barriers against every
reduction. A many-one reduction need not recover a satisfying assignment
from every small circuit: imposing that would demand an additional
witness-preservation property. The conditional obstruction to certain Levin
reductions recorded by the 2026 source concerns a gap problem and extra
assumptions, not an unconditional impossibility theorem for this exact
question. \sref{30006744}{1}

Membership in NP is unaffected by binary encoding of a huge bound $s$.
Every $N$-entry table has a decision-tree circuit of $O(N)$ gates. Bounds
above that universal ceiling are automatically yes; otherwise a circuit
certificate uses $O(N\log N)$ bits and is checked on all $N$ rows in
polynomial time. Neither this observation nor NP-hardness by itself
separates $\mathsf P$ from $\mathsf{NP}$.

The companion script checks the completion and edit identities on all
three-bit truth tables. These are finite Boolean checks, not circuit lower
bounds. Equation~\eqref{30006744:edit} and the counting proof apply at every
length without relying on the experiment.

\paragraph{Outcome.}
\textbf{Outcome: Exploratory approach with unresolved gap.}
The attempt isolates two quantifier failures in partial-to-total encodings
and proves an edit budget for a possible robust gadget. A polynomial-output
deterministic transformation whose unrestricted minimum circuit sizes
separate satisfiable from unsatisfiable instances is still missing. No such
hardness reduction or counterexample to its existence is established.

% END RESEARCH ADDITION 134
\subsection{Sources}
\begin{itemize}\raggedright
\item[\textnormal{[S1]}]\hypertarget{source:30006744:1}{} Halley Goldberg, Mandar Juvekar, and Valentine Kabanets, Non-Levin NP-Hardness of Implicit MCSP and PAC Learning under Few Assumptions, ECCC TR26-091, Introduction, printed p. 3 (2026).
\url{https://eccc.weizmann.ac.il/report/2026/091/download/}.
{\small\textit{Catalogue formulation source.}}
% BEGIN RESEARCH REFERENCES 134
\item[\textnormal{[E1]}]\hypertarget{extra:30006744:1}{} S. Hirahara and R. Ilango,
\emph{NP-hardness of the Minimum Circuit Size Problem from Well-Studied
Assumptions}, FOCS 2025, 1648--1664, abstract and introduction.
\url{https://www.rahulilango.com/papers/MCSP-Proceedings-2025.pdf}.
{\small\textit{Conditional quasipolynomial-time nonadaptive hardness;
not the target deterministic polynomial-time many-one reduction.}}

% END RESEARCH REFERENCES 134
\end{itemize}

\end{document}
